\documentclass[11pt]{article}
\usepackage[margin=1in]{geometry}
\usepackage{amsmath,amssymb,amsthm}
\usepackage{hyperref}
\usepackage{xurl}
\let\oldus\_
\renewcommand{\_}{\oldus\allowbreak}

\newtheorem{theorem}{Theorem}
\newtheorem{lemma}[theorem]{Lemma}
\theoremstyle{definition}
\newtheorem{definition}[theorem]{Definition}
\theoremstyle{remark}
\newtheorem{remark}[theorem]{Remark}

\title{A disproof of Conjecture 00000002348\\[2pt]
\large An expected condition number is at least $1$, never $m^{-1/2}$}
\date{}

\begin{document}
\sloppy
\maketitle

\section{The conjecture and its reading}

Conjecture 00000002348 reads: \emph{The expected condition number of the norm matrix of random
interpolation on smooth varieties is $m^{-1/2}$ ($m$ the sample count; a negative half power); the
expectation comes from the logarithmic-potential integral of the determinant of the norm matrix.}

The claim is a value for $\mathbb{E}[\kappa(A)]$, where $A$ is the (random) norm matrix built from
$m$ random samples. We refute two readings of ``is $m^{-1/2}$'':
\begin{itemize}
\item[(E)] \emph{exact:} $\mathbb{E}[\kappa(A_m)] = m^{-1/2}$ for the sample counts $m$;
\item[(O)] \emph{order of magnitude:} $\mathbb{E}[\kappa(A_m)] \le C\,m^{-1/2}$ for some constant $C$
and all large $m$. This is implied by $\mathbb{E}[\kappa(A_m)] = C m^{-1/2}$,
$\mathbb{E}[\kappa(A_m)] = \Theta(m^{-1/2})$, $\mathbb{E}[\kappa(A_m)] = O(m^{-1/2})$ and
$\mathbb{E}[\kappa(A_m)] \sim c\, m^{-1/2}$, so all of these are refuted too.
\end{itemize}
The construction of the norm matrix (the variety, the sampling law, the basis) is never used: the
theorems below hold for \emph{every} random matrix on \emph{every} probability space, with the matrix size
allowed to depend on $m$. Hence they apply to whatever norm matrix the conjecture intends. The remark
about the logarithmic-potential integral describes how the expectation would be computed and does
not change its value, so it is not used.

\paragraph{Conventions.}
Condition numbers take values in $[0,+\infty]$, with the standard convention $\kappa(A)=+\infty$ when
$A$ is singular (rank deficient). The expectation of a $[0,+\infty]$-valued random variable is its
Lebesgue integral. For measurable random variables this is the ordinary expectation; for arbitrary
$[0,+\infty]$-valued functions we use the lower Lebesgue integral (Mathlib's 	exttt{lintegral}), which
always exists, so no integrability or measurability assumption is needed for the bound. A
real-valued variant is in Theorem~\ref{thm:bochner}. Matrices are real and have at least one column
($n\ge1$, resp.\ $q\ge 1$): an empty $0\times 0$ matrix has $\|I\|=0$ and is excluded.
We use $m^{-1/2}=1/\sqrt m$ for $m\ge1$.

\section{Condition numbers are at least 1}

\begin{definition}[norm condition number]
Let $R$ be a nontrivial normed ring, i.e.\ a ring with a norm satisfying
$\|xy\|\le\|x\|\,\|y\|$ and $1\ne 0$. For $A\in R$ put $\kappa(A)=\|A\|\,\|A^{-1}\|$ if $A$ is invertible
and $\kappa(A)=+\infty$ otherwise. For real $n\times n$ matrices we use the spectral norm
$\|A\|_2=\sup_{\|x\|=1}\|Ax\|$ (Euclidean norms), which is submultiplicative; then
$\kappa(A)=\|A\|_2\|A^{-1}\|_2=\sigma_{\max}(A)/\sigma_{\min}(A)$ for invertible $A$.
\end{definition}

\begin{definition}[singular-value condition number]
For a real $p\times q$ matrix $A$ put $\sigma_{\max}(A)=\sup_{\|x\|=1}\|Ax\|$ and
$\sigma_{\min}(A)=\inf_{\|x\|=1}\|Ax\|$ (Euclidean norms, $x\in\mathbb{R}^q$). Define
$\kappa_2(A)=\sigma_{\max}(A)/\sigma_{\min}(A)$ if $\sigma_{\min}(A)>0$ and $\kappa_2(A)=+\infty$ if
$\sigma_{\min}(A)=0$, i.e.\ if $A$ does not have full column rank. This covers rectangular norm matrices,
for example an $m\times N$ matrix of evaluations at the samples.
\end{definition}

\begin{lemma}\label{lem:ge1}
(a) In a nontrivial normed ring, $\kappa(A)\ge 1$ for every $A$.
(b) For every real $p\times q$ matrix with $q\ge1$, $\kappa_2(A)\ge1$.
\end{lemma}

\begin{proof}
(a) If $A$ is singular, $\kappa(A)=+\infty$. Otherwise $\|1\|=\|1\cdot 1\|\le\|1\|^2$ and $\|1\|>0$ (as
$1\ne0$) give $\|1\|\ge1$, so $1\le\|1\|=\|AA^{-1}\|\le\|A\|\,\|A^{-1}\|$.
(b) If $\sigma_{\min}(A)=0$, $\kappa_2(A)=+\infty$. Otherwise pick the unit vector $e_1\in\mathbb{R}^q$:
$\sigma_{\min}(A)\le\|Ae_1\|\le\sigma_{\max}(A)$, and $\sigma_{\min}(A)$ is finite and positive, so
$\kappa_2(A)\ge\sigma_{\min}(A)/\sigma_{\min}(A)=1$.
\end{proof}

\section{The refutation}

\begin{lemma}\label{lem:exp}
If $(\Omega,P)$ is a probability space and $X:\Omega\to[0,+\infty]$ satisfies $X\ge1$ pointwise, then
$\mathbb{E}[X]=\int X\,dP\ge\int 1\,dP=1$.
\end{lemma}

\begin{theorem}[exact reading (E)]\label{thm:exact}
Let $m\ge2$, let $(\Omega,P)$ be any probability space, and let $A:\Omega\to\mathbb{R}^{n\times n}$
($n\ge1$) or $A:\Omega\to\mathbb{R}^{p\times q}$ ($q\ge1$) be any random matrix. Then
$\mathbb{E}[\kappa(A)]\ge1>m^{-1/2}$, resp.\ $\mathbb{E}[\kappa_2(A)]\ge1>m^{-1/2}$. In particular
neither expectation equals $m^{-1/2}$.
\end{theorem}

\begin{proof}
Lemmas~\ref{lem:ge1} and~\ref{lem:exp} give the lower bound $1$, and $m^{-1/2}<1$ because $m>1$ and
the exponent is negative.
\end{proof}

\begin{theorem}[order reading (O)]\label{thm:order}
Let $(\Omega_m,P_m)$ be probability spaces and $A_m$ random $N_m\times N_m$ matrices ($N_m\ge1$), resp.\
random $p_m\times q_m$ matrices ($q_m\ge1$), indexed by the sample count $m$. There are no $C\in\mathbb{R}$
and $M$ with $\mathbb{E}[\kappa(A_m)]\le C m^{-1/2}$ (resp.\ with $\kappa_2$) for all $m\ge M$.
\end{theorem}

\begin{proof}
Take $m=\max(M,\lceil C^2\rceil+1)$, so $m>C^2$. If $C\le0$ then $Cm^{-1/2}\le0<1$. If $C>0$ then
$C<\sqrt m$, so $Cm^{-1/2}=C/\sqrt m<1$. In both cases $\mathbb{E}\ge1>Cm^{-1/2}$.
\end{proof}

\begin{theorem}[real-valued expectation]\label{thm:bochner}
In any nontrivial normed ring, if $A$ is almost surely invertible and $\|A\|\,\|A^{-1}\|$ is integrable,
then $\int\|A\|\,\|A^{-1}\|\,dP\ge1$.
\end{theorem}

\begin{remark}[what is not refuted]
At $m=1$ we have $m^{-1/2}=1$, which is attainable (e.g.\ $A=I$), so (E) is refuted for every $m\ge 2$
only. We do not address non-literal rewordings that change the quantity, such as
``$\mathbb{E}[\kappa]-1$ is of order $m^{-1/2}$'' (a deviation-from-isometry statement) or
``$\mathbb{E}[\log\kappa]=m^{-1/2}$''. The conjecture states the value of the expected condition number
itself.
\end{remark}

\section{Relation to the earlier submission}

PR \#241 (orionsheep) gave the same argument ($\kappa\ge1$, hence $\mathbb{E}[\kappa]\ge1>m^{-1/2}$),
but it was closed without merging because its Lean capstone was the arithmetic fact
\texttt{(1:Nat) < 2 $\wedge$ 2 $\ne$ 1}. The reviewer wrote that ``the condition-number bound
$\kappa\ge1$ that constitutes the entire refutation of $E[\kappa]=m^{-1/2}$ appears only in prose, with no
matrices or expectations in Lean''. That PR also refuted only the exact value. The present Lean file
defines both condition numbers on actual real matrices and proves $\kappa\ge 1$ from submultiplicativity
and $\sigma_{\min}\le\sigma_{\max}$. It formalizes the expectation as a Lebesgue integral over an
arbitrary probability space and refutes both the exact reading at every $m\ge2$ and the order reading
for arbitrary $m$-indexed families.

\section{Lean formalization}

The file \texttt{lean/Conjecture2348/Basic.lean} (226 lines, Lean 4 / Mathlib \texttt{v4.33.1}, only
\texttt{import Mathlib}) contains, in namespace \texttt{C2348}:
\begin{itemize}
\item \texttt{condNum} (norm condition number in a normed ring, $+\infty$ for non-units) and
\texttt{one\_le\_condNum}; \texttt{condNum\_matrix\_eq} unfolds it, under the spectral norm
(\texttt{Matrix.Norms.L2Operator}), to $\|A\|_2\|A^{-1}\|_2$ with $A^{-1}$ the matrix inverse if
$\det A\ne0$, and to $+\infty$ if $\det A=0$;
\item \texttt{stretchMax}, \texttt{stretchMin}, \texttt{svCond} ($\sigma_{\max}/\sigma_{\min}$ via
\texttt{Matrix.toEuclideanLin}) and \texttt{one\_le\_svCond};
\item \texttt{one\_le\_expectation} (Lemma~\ref{lem:exp}, with \texttt{lintegral}),
\texttt{rpow\_neg\_half\_lt\_one}, \texttt{ne\_rpow\_neg\_half}, \texttt{not\_le\_C\_rpow\_neg\_half};
\item main theorems \texttt{conjecture\_2348\_false\_exact}, \texttt{conjecture\_2348\_false\_exact\_sv}
(Theorem~\ref{thm:exact}), \texttt{conjecture\_2348\_false\_order},
\texttt{conjecture\_2348\_false\_order\_sv} (Theorem~\ref{thm:order}) and
\texttt{one\_le\_integral\_condNum} (Theorem~\ref{thm:bochner}).
\end{itemize}
In the main theorems $m^{-1/2}$ is \texttt{(m:$\mathbb{R}$) \^{} (-(1/2:$\mathbb{R}$))}, embedded into
$[0,+\infty]$ by \texttt{ENNReal.ofReal}. \texttt{lake build} succeeds, and \texttt{\#print axioms} reports
only \texttt{propext}, \texttt{Classical.choice} and \texttt{Quot.sound} for every main theorem. There
is no \texttt{sorry} and no \texttt{native\_decide}.

\end{document}
